% problem_id: S001-proposition-annihilator-complex
% source_id: S001 (Stacks Project)
% source_file: local-cohomology.tex
% statement_locator: local-cohomology.tex lines 2776--2793
% proof_locator: local-cohomology.tex lines 2795--2947
% env: proposition
% status: text-extracted-verbatim (difficulty judgement pending, written by hand)
%
% The statement and proof below are the author's own text, extracted verbatim.
% Nothing is paraphrased, shortened, or supplemented.

\begin{proposition}
\label{proposition-annihilator-complex}
Let $A$ be a Noetherian ring which has a dualizing complex.
Let $T \subset T' \subset \Spec(A)$ be subsets stable under
specialization. Let $s \in \mathbf{Z}$. Let $K$ be an object of
$D_{\textit{Coh}}^+(A)$. The following are equivalent
\begin{enumerate}
\item there exists an ideal $J \subset A$ with $V(J) \subset T'$
such that $J$ annihilates $H^i_T(K)$ for $i \leq s$, and
\item for all $\mathfrak p \not \in T'$,
$\mathfrak q \in T$ with $\mathfrak p \subset \mathfrak q$
we have
$$
\text{depth}_{A_\mathfrak p}(K_\mathfrak p) +
\dim((A/\mathfrak p)_\mathfrak q) > s
$$
\end{enumerate}
\end{proposition}

\begin{proof}
This lemma is the natural generalization of
Proposition \ref{proposition-annihilator}
whose proof the reader should read first.
Let $\omega_A^\bullet$ be a dualizing complex. Let $\delta$ be its
dimension function, see Dualizing Complexes, Section
\ref{dualizing-section-dimension-function}.
An important role will be played by the finite $A$-modules
$$
E^i = \Ext_A^i(K, \omega_A^\bullet)
$$
For $\mathfrak p \subset A$ we will write $H^i_\mathfrak p$ to denote the
local cohomology of an object of $D(A_\mathfrak p)$ with respect to
$\mathfrak pA_\mathfrak p$. Then we see that
the $\mathfrak pA_\mathfrak p$-adic completion of
$$
(E^i)_\mathfrak p =
\Ext^{\delta(\mathfrak p) + i}_{A_\mathfrak p}(K_\mathfrak p,
(\omega_A^\bullet)_\mathfrak p[-\delta(\mathfrak p)])
$$
is Matlis dual to
$$
H^{-\delta(\mathfrak p) - i}_{\mathfrak p}(K_\mathfrak p)
$$
by
Dualizing Complexes, Lemma \ref{dualizing-lemma-special-case-local-duality}.
In particular we deduce from this the
following fact: an ideal $J \subset A$ annihilates
$(E^i)_\mathfrak p$ if and only if $J$ annihilates
$H^{-\delta(\mathfrak p) - i}_{\mathfrak p}(K_\mathfrak p)$.

\medskip\noindent
Set $T_n = \{\mathfrak p \in T \mid \delta(\mathfrak p) \leq n\}$.
As $\delta$ is a bounded function, we see that
$T_a = \emptyset$ for $a \ll 0$ and $T_b = T$ for $b \gg 0$.

\medskip\noindent
Assume (2). Let us prove the existence of $J$ as in (1).
We will use a double induction to do this. For $i \leq s$
consider the induction hypothesis $IH_i$:
$H^a_T(K)$ is annihilated by some $J \subset A$ with $V(J) \subset T'$
for $a \leq i$. The case $IH_i$ is trivial for $i$ small
enough because $K$ is bounded below.

\medskip\noindent
Induction step. Assume $IH_{i - 1}$ holds for some $i \leq s$.
Pick $J'$ with $V(J') \subset T'$ annihilating $H^a_T(K)$ for
$a \leq i - 1$ (the induction hypothesis guarantees we can
do this). We will show by descending induction on $n$
that there exists an ideal $J$ with $V(J) \subset T'$ such that the
associated primes of $J H^i_T(K)$ are in $T_n$.
For $n \ll 0$ this implies $JH^i_T(K) = 0$ 
(Algebra, Lemma \ref{algebra-lemma-ass-zero})
and hence $IH_i$ will hold.
The base case $n \gg 0$ is trivial because $T = T_n$ in this case
and all associated primes of $H^i_T(K)$ are in $T$.

\medskip\noindent
Thus we assume given $J$ with the property for $n$.
Let $\mathfrak q \in T_n$. Let $T_\mathfrak q \subset \Spec(A_\mathfrak q)$
be the inverse image of $T$. We have
$H^j_T(K)_\mathfrak q = H^j_{T_\mathfrak q}(K_\mathfrak q)$
by Lemma \ref{lemma-torsion-change-rings}.
Consider the spectral sequence
$$
H_\mathfrak q^p(H^q_{T_\mathfrak q}(K_\mathfrak q))
\Rightarrow
H^{p + q}_\mathfrak q(K_\mathfrak q)
$$
of Lemma \ref{lemma-local-cohomology-ss}.
Below we will find an ideal $J'' \subset A$ with $V(J'') \subset T'$
such that $H^i_\mathfrak q(K_\mathfrak q)$ is annihilated by $J''$ for all
$\mathfrak q \in T_n \setminus T_{n - 1}$.
Claim: $J (J')^i J''$ will work for $n - 1$.
Namely, let $\mathfrak q \in T_n \setminus T_{n - 1}$.
The spectral sequence above defines a filtration
$$
E_\infty^{0, i} = E_{i + 2}^{0, i} \subset \ldots \subset E_3^{0, i} \subset
E_2^{0, i} = H^0_\mathfrak q(H^i_{T_\mathfrak q}(K_\mathfrak q))
$$
The module $E_\infty^{0, i}$ is annihilated by $J''$.
The subquotients $E_j^{0, i}/E_{j + 1}^{0, i}$ for $i + 1 \geq j \geq 2$
are annihilated by $J'$ because the target of $d_j^{0, i}$
is a subquotient of
$$
H^j_\mathfrak q(H^{i - j + 1}_{T_\mathfrak q}(K_\mathfrak q)) =
H^j_\mathfrak q(H^{i - j + 1}_T(K)_\mathfrak q)
$$
and $H^{i - j + 1}_T(K)_\mathfrak q$ is annihilated by $J'$ by choice of $J'$.
Finally, by our choice of $J$ we have
$J H^i_T(K)_\mathfrak q \subset H^0_\mathfrak q(H^i_T(K)_\mathfrak q)$
since the non-closed points of $\Spec(A_\mathfrak q)$ have higher
$\delta$ values. Thus $\mathfrak q$ cannot be an associated prime of
$J(J')^iJ'' H^i_T(K)$ as desired.

\medskip\noindent
By our initial remarks we see that $J''$ should annihilate
$$
(E^{-\delta(\mathfrak q) - i})_\mathfrak q =
(E^{-n - i})_\mathfrak q
$$
for all $\mathfrak q \in T_n \setminus T_{n - 1}$.
But if $J''$ works for one $\mathfrak q$, then it works for all
$\mathfrak q$ in an open neighbourhood of $\mathfrak q$
as the modules $E^{-n - i}$ are finite.
Since every subset of $\Spec(A)$ is Noetherian with the induced
topology (Topology, Lemma \ref{topology-lemma-Noetherian}),
we conclude that it suffices
to prove the existence of $J''$ for one $\mathfrak q$.

\medskip\noindent
Since the ext modules are finite the existence of $J''$ is
equivalent to
$$
\text{Supp}(E^{-n - i}) \cap \Spec(A_\mathfrak q) \subset T'.
$$
This is equivalent to showing the localization of $E^{-n - i}$ at every
$\mathfrak p \subset \mathfrak q$, $\mathfrak p \not \in T'$
is zero. Using local duality over $A_\mathfrak p$ we find that we need
to prove that
$$
H^{i + n - \delta(\mathfrak p)}_\mathfrak p(K_\mathfrak p) =
H^{i - \dim((A/\mathfrak p)_\mathfrak q)}_\mathfrak p(K_\mathfrak p)
$$
is zero (this uses that $\delta$ is a dimension function).
This vanishes by the assumption in the lemma and $i \leq s$ and
our definition of depth in Definition \ref{definition-depth-complex}.

\medskip\noindent
To prove the converse implication we assume (2) does not hold
and we work backwards through the arguments above. First, we pick a
$\mathfrak q \in T$, $\mathfrak p \subset \mathfrak q$
with $\mathfrak p \not \in T'$ such that
$$
i = \text{depth}_{A_\mathfrak p}(K_\mathfrak p) +
\dim((A/\mathfrak p)_\mathfrak q) \leq s
$$
is minimal. Then
$H^{i - \dim((A/\mathfrak p)_\mathfrak q)}_\mathfrak p(K_\mathfrak p)$
is nonzero by the our definition of depth in
Definition \ref{definition-depth-complex}.
Set $n = \delta(\mathfrak q)$. Then
there does not exist an ideal $J \subset A$ with $V(J) \subset T'$
such that $J(E^{-n - i})_\mathfrak q = 0$.
Thus $H^i_\mathfrak q(K_\mathfrak q)$ is not
annihilated by an ideal $J \subset A$ with $V(J) \subset T'$.
By minimality of $i$ it follows from the spectral sequence displayed above
that the module $H^i_T(K)_\mathfrak q$
is not annihilated by an ideal $J \subset A$
with $V(J) \subset T'$. Thus $H^i_T(K)$
is not annihilated by an ideal $J \subset A$
with $V(J) \subset T'$. This finishes the proof of the proposition.
\end{proof}
